% EGA V, English translation by P. Blass, 1987 (cote 24): archive cote 24, site page 73 (PDF page 74).
% typescript. Draft compared: none.
% AI-transcribed by Claude Opus 5.5 (claude-opus-5-5, adaptive thinking, effort high), Message Batches API, batch ega24b (msgbatch_01RGJJeHdR96U16ZWFuUjtDx), 2026-11-13; not yet checked by eye.
\scanpage{24/73}
\mspage{73}

contains a k-rational point.  A point of $\underline{\mathrm{Grass}}_{n}(E)$ with values in S
i.e. a locally free quotient module G of rank n of E canonically
defines a subscheme of P(E), (i.e., to say) P(G).  Such a
subscheme (but without \uline{imposing} the rank of G) is called a \uline{linear}
\uline{subvariety} of P(E) (relative to S if there is a possibility of
confusion).  It is therefore a projective fibration of relative
dimension (n-1) if $n \geq 1$, (and empty is $n = 0$).  We immediately
verify that the section of $\underline{\mathrm{Grass}}_{n}(E)$ i.e. G is known if we know the
linear subvariety corresponding to P(E).  In this manner the
grassmanian can be interpreted as representing the functor
(``linear subvarieties of relative dimension n - 1 of $\underline{P}_{S}{}^{1}$ '') for $S^{1}$
variable $n \geq 1$.  It is furthermore possible to give an intrinsic
characterization of the latter functor i.e. of the notion of
linear subvariety of relative dimension m and that are of
``projective degree one'' at every $s \in S$; this characterization will
be given in a later chapter and we shall not need it at all here.

Let us again suppose that E is locally free of rank r, let $E^{\vee}$
be its dual.  Then by a polarity we find a canonical isomophism
$\underline{\mathrm{Grass}}_{n}(E) \simeq \underline{\mathrm{Grass}}_{r-n}(E^{\vee})$ that assigns to a quotient $\underline{G}$ of $\underline{E}$ the
quotient $\underline{E}^{\vee}/\underline{G}^{\vee}$ of $\underline{E}^{\vee}$.  From the point of view of linear varieties
to a linear variety L of relative dimension m of P there
corresponds the linear dual variety $L^{\circ}$ of relative dimension (r-1)
- 1 - m of $P^{\vee}$ i.e. of relative codimension (m+1) in $P^{\vee}$.  (N.B. (r-%
1) is here the relative common dimension of P and $P^{\vee}$ over S), which
we may visualize geometrically \uline{as follows}.  Let us first of all
take $n = r - 1$, we find an isomorphism $\underline{P}(E^{\vee}) \simeq \underline{\mathrm{Grass}}_{r-1}(\underline{E})$ that
allows us to identify the points of $P^{\vee}$ with values in S (let us
